\chapter{概率工具与集中不等式}

随机算法分析不是把“随机”留给实验，而是把随机选择写成随机变量，再用期望、方差和尾界控制失败事件。本章建立后续反复使用的工具箱。

\section{事件、指示变量与期望线性性}

对事件 $A$，定义指示变量
\[
\ind_A=
\begin{cases}
1,&A\text{ 发生},\\
0,&A\text{ 不发生}.
\end{cases}
\]
它满足 $\E[\ind_A]=\Prb(A)$。计数型随机变量常可写成指示变量之和，从而利用
\[
\E\left[\sum_iX_i\right]=\sum_i\E[X_i].
\]
期望线性性不要求 $X_i$ 独立。

\begin{example}[空桶数]
把 $n$ 个球独立均匀放入 $n$ 个桶。令 $I_j$ 表示第 $j$ 个桶为空，则
\[
\Prb(I_j=1)=\left(1-\frac1n\right)^n.
\]
空桶总数 $X=\sum_{j=1}^nI_j$，所以
\[
\E X=n\left(1-\frac1n\right)^n\approx\frac ne.
\]
各桶是否为空并不独立，但不影响期望计算。
\end{example}

\section{Union bound}

\begin{theorem}[Union bound]
对任意事件 $E_1,\dots,E_m$，
\[
\Prb\left(\bigcup_{i=1}^mE_i\right)
\le\sum_{i=1}^m\Prb(E_i).
\]
\end{theorem}

它不要求事件独立。代价是可能重复计算交叠部分，因此界可能较松。算法中常把“存在某一步失败”分解成多个局部失败事件，再控制总失败概率。

\begin{example}[对所有查询同时成功]
若单个查询失败概率至多 $\delta/q$，则对预先固定的 $q$ 个查询，至少一个失败的概率至多 $\delta$。若查询是根据前面结果自适应选择的，还需检查单查询保证是否在条件分布下仍成立，不能机械套用。
\end{example}

\section{Markov 与 Chebyshev 不等式}

\begin{theorem}[Markov 不等式]
若 $X\ge0$ 且 $a>0$，则
\[
\Prb[X\ge a]\le\frac{\E X}{a}.
\]
\end{theorem}

\begin{proof}
点态上有 $X\ge a\ind_{\{X\ge a\}}$。两边取期望：
\[
\E X\ge a\Prb[X\ge a].
\]
\end{proof}

\begin{theorem}[Chebyshev 不等式]
若 $X$ 方差有限，则对 $t>0$，
\[
\Prb[|X-\E X|\ge t]\le\frac{\Var(X)}{t^2}.
\]
\end{theorem}

\begin{proof}
对非负随机变量 $(X-\E X)^2$ 使用 Markov 不等式，阈值取 $t^2$。
\end{proof}

Markov 只使用均值，Chebyshev 使用到二阶矩。它们假设少、适用广，但通常只给多项式衰减。

\section{矩母函数与 Chernoff 方法}

对任意 $\lambda>0$，指数函数单调，故
\begin{align*}
\Prb[X\ge a]
&=\Prb[e^{\lambda X}\ge e^{\lambda a}]\\
&\le e^{-\lambda a}\E[e^{\lambda X}].
\end{align*}
随后选择最优 $\lambda$。若 $X=\sum_iX_i$ 且 $X_i$ 独立，则
\[
\E[e^{\lambda X}]=\prod_i\E[e^{\lambda X_i}],
\]
这正是独立性带来指数尾界的关键位置。

\begin{theorem}[Bernoulli 和的 Chernoff 界]
设 $X_i\in\{0,1\}$ 相互独立，$X=\sum_iX_i$，$\mu=\E X$。则对 $\delta>0$，
\[
\Prb[X\ge(1+\delta)\mu]
\le\left(\frac{e^\delta}{(1+\delta)^{1+\delta}}\right)^\mu.
\]
特别地，对 $0<\delta\le1$，
\[
\Prb[X\ge(1+\delta)\mu]
\le e^{-\mu\delta^2/3},
\]
并且
\[
\Prb[X\le(1-\delta)\mu]
\le e^{-\mu\delta^2/2}.
\]
\end{theorem}

\begin{proof}[上尾证明骨架]
对 Bernoulli 变量，
\[
\E[e^{\lambda X_i}]
=1-p_i+p_ie^\lambda
\le\exp(p_i(e^\lambda-1)).
\]
独立性给出
\[
\E[e^{\lambda X}]\le\exp(\mu(e^\lambda-1)).
\]
代入 Markov 型界并令 $e^\lambda=1+\delta$，得到第一条。再用
$(1+\delta)\ln(1+\delta)-\delta\ge\delta^2/3$（$0<\delta\le1$）得到简化形式。
\end{proof}

\begin{warningbox}
Chernoff 界不是“任何随机变量都指数集中”。标准版本需要独立 Bernoulli 或相应负相关条件。若算法反复复用同一随机哈希，变量可能相关，必须重新检查假设。
\end{warningbox}

\section{Hoeffding 不等式}

\begin{theorem}[Hoeffding]
若独立随机变量 $X_i\in[a_i,b_i]$，令 $S=\sum_iX_i$，则
\[
\Prb[S-\E S\ge t]
\le\exp\left(-\frac{2t^2}{\sum_i(b_i-a_i)^2}\right).
\]
双侧界多一个因子 2。
\end{theorem}

对 $n$ 个 $[0,1]$ 样本的均值 $\bar X$，
\[
\Prb[|\bar X-\E X|\ge\epsilon]
\le2e^{-2n\epsilon^2}.
\]
要使失败概率至多 $\delta$，充分取
\[
n\ge\frac{1}{2\epsilon^2}\ln\frac2\delta.
\]

\section{概率放大}

若一次 Monte Carlo 判定以概率至少 $2/3$ 正确，独立运行 $k$ 次并取多数票，错误概率由 Chernoff 界降到 $e^{-\Theta(k)}$。因此要达到错误概率 $\delta$，通常只需 $k=O(\log(1/\delta))$ 次。

对估计问题，常使用 median trick：每组先产生常数成功概率的估计，再对多组独立估计取中位数。只要超过半数组合格，中位数就合格。

\begin{intuitionbox}
随机算法中常把任务拆成两层：先构造一个常数成功概率、容易分析的基本估计器；再用独立重复把可靠性提升到用户指定的 $1-\delta$。
\end{intuitionbox}

\section{工具选择表}

\begin{center}
\begin{tabularx}{\textwidth}{lXXX}
\toprule
工具 & 所需信息 & 典型尾部 & 常见用途\\
\midrule
Union bound & 单个事件概率 & 概率相加 & 多个失败事件\\
Markov & 非负、均值 & $1/a$ & 粗略单侧界\\
Chebyshev & 均值、方差 & $1/t^2$ & 无强独立假设估计\\
Chernoff & 独立 Bernoulli 和 & 指数衰减 & 计数、采样、放大\\
Hoeffding & 独立有界变量 & 指数衰减 & 样本均值\\
\bottomrule
\end{tabularx}
\end{center}

\section{小结与验收}

\begin{checkpointbox}
\begin{enumerate}
  \item 用指示变量重新推导空桶数期望；
  \item 从 Markov 推导 Chebyshev；
  \item 在 Chernoff 证明中圈出使用独立性的那一步；
  \item 推导估计 Bernoulli 均值到加性误差 $\epsilon$、失败概率 $\delta$ 所需的样本量；
  \item 解释为什么重复实验必须使用独立随机性才能直接应用多数票放大。
\end{enumerate}
\end{checkpointbox}
